\documentclass[11pt]{article}
\usepackage{mystyle}

\title{Rosen --- \S 9.1: Relations and Their Properties\\ \large Selected Exercises}
\author{Alston}
\date{Semester 2, 2026}

\begin{document}
\maketitle

% ======================================================================
\begin{exercise}{8}
    Show that the relation $R = \emptyset$ on a nonempty set $S$ is
    symmetric and transitive, but not reflexive.
\end{exercise}

% ======================================================================
\begin{exercise}{10}
    Give an example of a relation on a set that is
    \begin{enumerate}[label=(\alph*)]
        \item both symmetric and antisymmetric.
        \item neither symmetric nor antisymmetric.
    \end{enumerate}
\end{exercise}

% ======================================================================
\begin{exercise}{15}
    Can a relation on a set be neither reflexive nor irreflexive?
\end{exercise}

% ======================================================================
\begin{exercise}{22}
    \textit{A relation $R$ is called \emph{asymmetric} if
    $(a,b) \in R$ implies that $(b,a) \notin R$.}

    Must an asymmetric relation also be antisymmetric? Must an
    antisymmetric relation be asymmetric? Give reasons for your
    answers.
\end{exercise}

% ======================================================================
\begin{exercise}{25}
    How many different relations are there from a set with $m$
    elements to a set with $n$ elements?
\end{exercise}

% ======================================================================
\begin{exercise}{27}
    \textit{Let $R$ be a relation from a set $A$ to a set $B$. The
    inverse relation from $B$ to $A$, denoted by $R^{-1}$, is the set
    of ordered pairs $\{(b,a) \mid (a,b) \in R\}$. The complementary
    relation $\overline{R}$ is the set of ordered pairs
    $\{(a,b) \mid (a,b) \notin R\}$.}

    Let $R$ be the relation $R = \{(a,b) \mid a \text{ divides } b\}$
    on the set of positive integers. Find
    \begin{enumerate}[label=(\alph*)]
        \item $R^{-1}$.
        \item $\overline{R}$.
    \end{enumerate}
\end{exercise}

% ======================================================================
\begin{exercise}{29}
    Suppose that the function $f$ from $A$ to $B$ is a one-to-one
    correspondence. Let $R$ be the relation that equals the graph of
    $f$. That is, $R = \{(a, f(a)) \mid a \in A\}$. What is the
    inverse relation $R^{-1}$?
\end{exercise}

% ======================================================================
\begin{exercise}{33}
    Let $R$ be the relation on the set of people consisting of pairs
    $(a,b)$, where $a$ is a parent of $b$. Let $S$ be the relation on
    the set of people consisting of pairs $(a,b)$, where $a$ and $b$
    are siblings (brothers or sisters). What are $S \circ R$ and
    $R \circ S$?
\end{exercise}

% ======================================================================
\begin{exercise}{39}
    Let $R$ be the relation on the set of people with doctorates such
    that $(a,b) \in R$ if and only if $a$ was the thesis advisor of
    $b$. When is an ordered pair $(a,b)$ in $R^2$? When is an ordered
    pair $(a,b)$ in $R^n$, when $n$ is a positive integer? (Assume
    that every person with a doctorate has a thesis advisor.)
\end{exercise}


% ======================================================================
\begin{exercise}{47}
    How many relations are there on a set with $n$ elements that are
    \begin{enumerate}[label=(\alph*)]
        \item symmetric?
        \item antisymmetric?
        \item asymmetric?
        \item irreflexive?
        \item reflexive and symmetric?
        \item neither reflexive nor irreflexive?
    \end{enumerate}
\end{exercise}

% ======================================================================
\begin{exercise}{48}
    How many transitive relations are there on a set with $n$
    elements if
    \begin{enumerate}[label=(\alph*)]
        \item $n=1$?
        \item $n=2$?
        \item $n=3$?
    \end{enumerate}
\end{exercise}

% ======================================================================
\begin{exercise}{49}
    Find the error in the ``proof'' of the following ``theorem.''

    \textit{``Theorem'': Let $R$ be a relation on a set $A$ that is
    symmetric and transitive. Then $R$ is reflexive.}

    \textit{``Proof'': Let $a \in A$. Take an element $b \in A$ such
    that $(a,b) \in R$. Because $R$ is symmetric, we also have
    $(b,a) \in R$. Now using the transitive property, we can
    conclude that $(a,a) \in R$ because $(a,b) \in R$ and
    $(b,a) \in R$.}
\end{exercise}

% ======================================================================
\begin{exercise}{52}
    Show that the relation $R$ on a set $A$ is antisymmetric if and
    only if $R \cap R^{-1}$ is a subset of the diagonal relation
    $\Delta = \{(a,a) \mid a \in A\}$.
\end{exercise}

% ======================================================================
\begin{exercise}{55}
    Let $R$ be a relation that is reflexive and transitive. Prove
    that $R^n = R$ for all positive integers $n$.
\end{exercise}

% ======================================================================
\begin{exercise}{58}
    Let $R$ be a symmetric relation. Show that $R^n$ is symmetric for
    all positive integers $n$.
\end{exercise}

% ======================================================================
\begin{exercise}{59}
    Suppose that the relation $R$ is irreflexive. Is $R^2$
    necessarily irreflexive? Give a reason for your answer.
\end{exercise}

\end{document}